% Transcription — fonds Alexandre Grothendieck, shelfmark 126, batch 2 (pages 21-34).
% Established from the facsimile archives/batches/126/batch-02.pdf, prepared from
% a copy of 126.pdf supplied by the user (PDF page k = archivists' page 20+k,
% no cover sheet), per the skill transcribe-grothendieck. The folder is
% thirty-four pages; this is the second and last batch (pages 1-20 are batch 1,
% transcribed separately and in parallel, so the notation below was fixed from
% these pages alone).
%
% Pass: Opus 5.5 (claude-opus-5-5), 2026-09-24 — first pass, unchecked against the pages by a human.
%
% Pages skipped: 24, a blank sheet of the typist's letterhead (M. Dufour,
% varitypiste) with nothing in his hand; 25 and 30, blank sheets (only
% show-through). No page is summarised.
%
% Pages transcribed: 21-23, 26-29, 31-34, all from his hand. Pages 21-23 are
% blue ink on the typist's letterhead paper (the letterhead shows through
% from the verso) and are very fast drafting: the formulas and the lemma
% statements carry, most of the connective prose does not and is left \ill{}.
% Pages 26-29 and 31-34 are black ink and somewhat steadier. He does not
% paginate; no internal cross-reference to a page number occurs.
%
% What the batch is about. Two runs on descent.
% (1) Pages 21-23 (the start of the run is presumably at the end of batch 1):
% descent of étale coverings along a finite radicial-type surjection Y' -> Y,
% by successive extensions; a « Lemme 2 » (a finite surjective Y' -> Y over
% which X' = X x_Y Y' splits into components of (separable?) degree 1 over Y)
% and a « Lemme 3 » (an étale surjective Y* -> Y with the same property),
% the trivial covering Y* x I built on the set I of connected components,
% and the descent data compared on it; the page ends with « À formuler... »
% about finite universally open morphisms with connected (?) fibres, and a
% question: characterising functors representable by étale coverings.
% (2) Pages 26-29: descent « generically ». S noetherian, T/S of finite type
% and flat, U dense constructible in T x_S T, V dense constructible in the
% triple product, F a fibred category, X in F_S, and phi an isomorphism
% (X x_S T)|U -> (T x_S X)|U satisfying the cocycle condition over V.
% Proposition: there is a dense constructible T' in T and a genuine descent
% datum phi' on X x_S T' agreeing with phi over U cap (T' x_S T'). Proof by
% shrinking U, V, T; W = p12^{-1}(U) cap p23^{-1}(U) with pi = p13 : W -> T x T;
% psi : pi^*(X_1) -> pi^*(X_2) as the composite through the middle factor;
% compatibility with pi checked through the five-point scheme W' (the
% diagram of page 27) and the equality phi_{t2 t'} phi_{t' t1} =
% phi_{t2 t''} phi_{t'' t1}; phi' obtained because W -> T x_S T is a morphism
% of F-descent; cocycle for phi' through a six-point scheme W''' (diagram of
% page 28). A « Condition de représentation » closes it. (3) Pages 31-34: the
% question of descent along a faithfully flat morphism that is not finite.
% Typical case: T finite étale Galois over S with group G, s a closed point,
% T' = T - (f^{-1}(s) - (t)). T x_S T is a disjoint union of the graphs U_g of
% the elements of G; T' x_S T' is the diagonal plus copies of T'' =
% T - f^{-1}(s); the same for the triple product with V_{g,g'}. Conclusion
% (underlined): a descent datum on X over T' -> S is the same thing as an
% action of G on X|T'' compatible with that on T''; effectivity reduces to
% extending X_0 over S - s across s.
%
% Words he underlines are set in \emph (the reading view has no text
% underline). Notation fixed once for the batch: his script U is \mathcal{U} (on pages 26-
% 29 and in the diagrams, where the arrows are labelled with it), V and W
% roman; p^{(3)}_{ij} for the projections of the triple product (he writes
% the (3) as a superscript, the -1 after it); \varphi_{t_2 t'} and the like
% for the components of phi at points, as he writes them; the asterisks of
% X^*, Y^*, X'^* as written; \Delta_T and \Delta_3 T as written on
% pages 31-32; X|T'' for his restriction bar.
%
% Readings to check first: the whole prose of pages 21-23 (the ink is pale
% and the hand very fast); the numbering of the two lemmas on pages 21-22;
% « À formuler » on page 23; the heavily corrected line « F catégorie
% fibrée ... » on page 26; the tuple lists defining W' (page 27) and W'''
% (page 28); the last words of page 28 (a name?); « Condition de
% représentation » on page 29; the formula T x_S T = Delta_T ... (G x G) on
% page 31, where a symbol before (G x G) is struck.
%
% The dating is the inventory's, copied verbatim: the folder itself is s.d.,
% and its inventory group « Descentes » carries the date quoted. The title in
% \foldertitle is the archivists' (not bracketed in the inventory).

\documentclass[11pt,a4paper]{article}
% Le préambule commun à toutes les transcriptions du fonds Grothendieck.
%
% Il tient en une idée : **l'incertitude se compose, elle ne se résout pas.**
% Une transcription qui lisse un mot douteux a détruit la seule chose pour
% laquelle on la faisait — un lecteur doit pouvoir savoir, à chaque endroit,
% ce qui était sur le feuillet et ce qui a été ajouté. Les six macros qui
% suivent sont donc l'appareil critique, et non de la décoration.
%
% Les mêmes commandes sont reconnues par `scripts/render.mjs`, qui produit la
% vue de lecture du site. Une macro ajoutée ici doit l'être aussi là-bas,
% faute de quoi le rendu échouera bruyamment — ce qui est voulu : un rendu
% silencieusement faux est pire qu'un rendu absent.

\ProvidesPackage{grothendieck}[2026/08/08 Transcriptions du fonds Grothendieck]

\usepackage{amsmath,amssymb,amsthm}
\usepackage{mathrsfs}
\usepackage{stmaryrd}
% Les diagrammes commutatifs sont partout dans ce fonds : c'est la langue
% même de la géométrie algébrique de Grothendieck, et une transcription qui
% les rendrait en prose ne transcrirait rien.
\usepackage{tikz-cd}
% Français par défaut — c'est la langue des feuillets — mais l'anglais est
% chargé aussi : la traduction et le résumé sont des documents anglais, et la
% typographie française (espaces avant les deux-points) y serait une faute.
% Ces documents basculent par \selectlanguage{english} en tête de corps.
\usepackage[english,french]{babel}
\usepackage{fontspec}
\usepackage{geometry}
\geometry{a4paper,margin=25mm,marginparwidth=28mm}
\usepackage{marginnote}
\usepackage{xcolor}
% ulem, not soul. `soul`'s \st and \ul are built on a character-by-character
% scan that XeLaTeX's font machinery breaks: the document fails with an
% undefined control sequence rather than degrading. ulem does the same two jobs
% and is Unicode-engine safe. `normalem` stops it hijacking \emph, which the
% transcriptions use for Grothendieck's own emphasis.
\usepackage[normalem]{ulem}
\usepackage{hyperref}
\hypersetup{colorlinks=true,linkcolor=black,urlcolor=black,pdfborder={0 0 0}}

% --- Métadonnées du lot -----------------------------------------------------
% Reprises telles quelles par le script de rendu, qui les affiche en tête de
% la vue de lecture. Elles fixent ce que le fichier prétend couvrir : sans
% elles, une transcription flotte sans point d'attache dans un fonds de seize
% mille feuillets.

\newcommand{\folder}[1]{\gdef\grothFolder{#1}}
\newcommand{\batch}[1]{\gdef\grothBatch{#1}}
\newcommand{\leaves}[2]{\gdef\grothFirst{#1}\gdef\grothLast{#2}}
\newcommand{\foldertitle}[1]{\gdef\grothFoldertitle{#1}}
\gdef\grothFolder{?}\gdef\grothBatch{?}\gdef\grothFirst{?}\gdef\grothLast{?}\gdef\grothFoldertitle{}

% --- Le résumé --------------------------------------------------------------
%
% Il ouvre la lecture modernisée, sous le titre « Résumé », et il est composé
% en petit. Ce n'est pas un ornement : c'est le seul passage du document qui
% ne soit pas des mathématiques, et un lecteur doit pouvoir le reconnaître
% avant de l'avoir lu — puis le sauter s'il connaît déjà le sujet, ou s'y
% arrêter s'il ne le connaît pas. Le filet qui le suit marque la reprise du
% corps.

\newenvironment{resume}
  {\small\setlength{\parskip}{0.45em}}
  {\par\medskip\hrule\medskip}

% --- Mots-clés ---------------------------------------------------------------
%
% En anglais, en clôture du résumé de la lecture modernisée : le vocabulaire
% moderne sous lequel chercher ce que les feuillets construisent. Le manifeste
% du site (`npm run manifest`) les extrait pour étiqueter la cote entière —
% cette ligne est la source unique des tags, il n'y a pas de fichier à côté.
\newcommand{\keywords}[1]{\par\smallskip\noindent
  {\footnotesize\textsc{Keywords} --- \textit{#1}}\par}

% --- L'appareil critique ----------------------------------------------------

% Le numéro de feuillet, en marge. C'est l'ancre qui permet de régler un
% désaccord sur une formule en regardant *un* feuillet plutôt qu'une chemise
% de six cents. Le site s'en sert aussi pour tourner les pages du fac-similé
% quand on fait défiler la transcription.
\newcommand{\leaf}[1]{%
  \marginnote{\footnotesize\color{gray!70}\textsf{#1}}%
  \label{leaf:#1}\ignorespaces}

% Illisible : on ne devine pas. Un mot inventé qui se lit comme les autres est
% le pire résultat possible.
\newcommand{\ill}{\textcolor{red!60!black}{[\,\ldots\,]}}

% Lecture douteuse : le mot est donné, et signalé comme incertain.
\newcommand{\uncertain}[1]{\uline{#1}}

% Ajout de l'éditeur — une lettre manquante, un mot sous-entendu.
\newcommand{\add}[1]{\textcolor{blue!50!black}{[#1]}}

% Biffé par Grothendieck lui-même. Ce qu'il a rayé fait partie du document :
% une rature dit souvent où la pensée a changé de direction.
\newcommand{\struck}[1]{\sout{#1}}

% Note du transcripteur, sur l'état matériel du feuillet.
\newcommand{\note}[1]{\par\smallskip{\small\color{gray!60!black}\textit{[#1]}}\par\smallskip}

% Note marginale de Grothendieck, distincte du corps du texte.
\newcommand{\marginal}[1]{\par\smallskip\noindent\hspace{1em}%
  {\small\color{gray!60!black}#1}\par\smallskip}

% --- Titre ------------------------------------------------------------------

\AtBeginDocument{%
  \begin{center}
    {\large\bfseries Fonds Alexandre Grothendieck --- cote n\textsuperscript{o}~\grothFolder}\\[2pt]
    {\itshape\grothFoldertitle}\\[4pt]
    {\small feuillets \grothFirst--\grothLast\ (lot \grothBatch)}\\[2pt]
    {\footnotesize\color{gray!60!black}Transcription établie d'après le fac-similé numérisé
     par l'Université de Montpellier.\\
     Les lectures incertaines sont soulignées, les illisibles notées [\,\ldots\,],
     les ajouts entre crochets.}
  \end{center}
  \vspace{1em}}

\endinput


\folder{126}
\batch{2}
\pages{21}{34}
\dating{s.d. — le groupe « Descentes » (126 à 133) est daté [avant 1970]}
\watermark{Édition de démonstration}
\foldertitle{Descente fidèlement plate. Divers : notes manuscrites (s.d.).}

\begin{document}

\section*{Revêtements étales et descente}

\note{les pages 21 à 23 sont d'une écriture très rapide, à l'encre bleue ;
les formules et les énoncés passent, la plus grande partie de la prose de
liaison non}

\page{21}

Faisons \uncertain{une} \uncertain{extension}, de la \ill{} \ill{}
\uncertain{successives} $Y_1 = X$, [d'où $X_1^{*} = X \times_Y Y_1$,
$X_1 = X_{(1)} \amalg X^{(1)}$, puis $Y_2 = X^{(1)}$] \ill{} \ill{}

\uncertain{Lemme} 2. Il existe un $Y'_2$ \uncertain{fini} sur $Y$, avec
$Y'_2 \to Y$ surjectif [et \struck{\uncertain{radiciel}}
\uncertain{universellement ouvert} \ill{} \ill{}] tel que
$X' = X \times_Y Y'$ se \uncertain{décompose} \ill{} \ill{} \ill{}
\uncertain{réunion} $X'_i$ \uncertain{connexes} de degré
\struck{\uncertain{fini}} \uncertain{séparable} 1 sur $Y$.
\note{énoncé encadré d'un trait vertical à gauche ; l'indice de $Y'_2$ est
écrit, puis omis dans $X' = X \times_Y Y'$}

Soit $I$ l'ensemble des composantes connexes de $X'$, considérons
$Y' \times I$ comme \ill{} rev.\ étale de $Y'$. Il en \uncertain{résulte}
\ill{} \ill{} \ill{} $X'_0 = Y' \times I$ \ill{} \ill{} \ill{} de descente
\uncertain{relatives} : $Y' \to Y$, d'où un \emph{rev.}\
\emph{étale} $X$ de $Y$, \ill{} $X'^{*}_0$ \uncertain{par}
\uncertain{image} \uncertain{inverse}. Soit $Y^{*}_{\ill{}}$ un
\uncertain{rev.} \emph{étale} [\uncertain{disons} \ill{}] de $Y$ qui
\uncertain{déploie} $X_0$ complètement, \uncertain{puis} l'\ill{} \ill{}
de la situation \uncertain{sur}

\note{schéma dans la marge gauche :}

\begin{tikzcd}[column sep=small, row sep=small]
X_0^{*} \arrow[d] & X'^{*}_0 \arrow[l] \arrow[d] \\
Y & Y' \arrow[l] \\
Y^{*}_{\ill{}} \arrow[u] & Y'^{*}_{\ill{}} \arrow[l] \arrow[u]
\end{tikzcd}

\note{les flèches verticales du haut descendent de $X^{*}_0$, $X'^{*}_0$
vers $Y$, $Y'$ ; celles du bas montent vers $Y$, $Y'$. Les indices sous
$Y^{*}$ et $Y'^{*}$ sont surchargés ; la lettre $Y$ du bas à gauche est
repassée}

\page{22}

$Y^{*}_{\ill{}}$. Donc maintenant \ill{} \ill{} \ill{} \ill{} \ill{} la
\ill{} de descente de $X'^{*}$ donne un $Y'^{*}$ complètement
\uncertain{décomposé}. Il en résulte aisément que les composantes connexes
de $X'^{*}$ sont les images inverses des composantes connexes de $X^{*}$
(\uncertain{cf.}\ \emph{\uncertain{lemme}} 3) compte tenu
\uncertain{que} \struck{\ill{}} $Y'^{*} \to Y^{*}$ est
\uncertain{radiciel} \uncertain{surjectif}.

Donc

\emph{\uncertain{Lemme}} \uncertain{3}. Il existe un rev.\ étale
surjectif $Y^{*} \to Y$ tel que $X^{*} = X \times_Y Y^{*}$ ait des
composantes connexes qui \uncertain{soient} de degré 1 sur $Y^{*}$.
\note{énoncé encadré d'un double trait vertical}

Le morphisme canonique
\[
X^{*} \longrightarrow Y^{*} \times I = X^{*}_0
\]
[$I$ \ill{} \ill{} \ill{} \uncertain{composantes} \uncertain{connexes} de
$X^{*}$] est compatible \struck{\ill{}} \uncertain{avec} les
\uncertain{données} de descente sur $Y^{*} \times I$, (et $X^{*}$),
\ill{}

\page{23}

\uncertain{donc} \uncertain{définit} un \uncertain{morphisme}
\struck{\ill{}}
\[
X \longrightarrow X_0
\]
\ill{} \ill{} \uncertain{surjectif} \uncertain{radiciel} \uncertain{fini}.
\uncertain{Pour} \uncertain{montrer} que \ill{} \uncertain{morphisme} est
bien \ill{}, \ill{} \uncertain{universel} que \ill{} \ill{} \ill{}
\uncertain{dit}, \uncertain{on} \uncertain{voit} que \ill{} $X^{*}$ \ill{}
\ill{} $Y^{*}$, un \uncertain{morphisme} \ill{} \ill{} \ill{} $X'^{*}$ et
$Y^{*}$ \ill{} \ill{} \ill{} \ill{}
\[
X^{*} \longrightarrow X^{*}_0 \longrightarrow Y'^{*}
\]
[et \uncertain{si} $Y'^{*}$ \uncertain{provient} d'un $Y'$
\uncertain{morphisme} \struck{\ill{}} \ill{} \ill{} \ill{}
\struck{\ill{}} \ill{} \ill{} \ill{}, $X^{*}_0 \to Y'^{*}$ \ill{} \ill{}].
\uncertain{Pour} \uncertain{vérifier} \ill{} \ill{} \ill{} \ill{} \ill{}
\ill{}, \ill{} \ill{} \ill{} $X^{*}_0 = Y^{*}$, \ill{} \ill{}
\uncertain{immédiat} \ill{} (\ill{} \ill{}).

\emph{À \struck{\ill{}} \uncertain{formuler}} \ill{} \ill{} \ill{}
\ill{} \ill{} : \ill{} \struck{\ill{}} \emph{\uncertain{univ.}!\
\uncertain{ouvert}} $X \to Y$ \ill{} \ill{} \uncertain{connexité}
\uncertain{géométrique} \uncertain{des} \uncertain{fibres}
(\uncertain{ouvert} \uncertain{réduit} \uncertain{fini} \ill{}),
« \ill{} \ill{} \ill{} » i.e.\ \ill{} \ill{} \uncertain{ouvert}
\struck{\ill{}}

\emph{\uncertain{Question}} \uncertain{Caractérisation} des foncteurs
représentables par des rev.\ étales ???

\section*{Descente générique}

\page{26}

$S$ schéma noethérien

$T/S$ de type fini, \struck{\ill{}} plat

$\mathcal{U}$ partie \uncertain{constructible} dense de $T \times_S T$,
$V$ partie \uncertain{constructible} dense de $T \times_S T \times_S T$,
$V \subset \bigcap_{i \neq j} p^{(3)-1}_{ij}(\mathcal{U})$

$F$ catégorie fibrée sur la catégorie des $S$-schémas de type fini, on
suppose \struck{\ill{}} \uncertain{les} morphismes \struck{\uncertain{fid.\
plats} et \ill{} \ill{} \ill{} \ill{}} \ill{} \struck{\uncertain{$F$-descente}}
\struck{\ill{}}
\note{ligne très surchargée : plusieurs mots biffés, un ajout
interlinéaire « \uncertain{$S$-schémas de type fini} » et un second
passage biffé d'un double trait ; on ne donne que ce qui se lit}

$X \in F_S$
\[
\varphi \colon (X \times_S T)|\mathcal{U} \xrightarrow{\ \sim\ }
(T \times_S X)|\mathcal{U}
\]
on suppose $p_{13}^{*}(\varphi) = p_{23}^{*}(\varphi)\, p_{12}^{*}(\varphi)$
au-dessus de $V$.

\emph{Proposition}. Sous ces conditions, il existe une partie
\uncertain{constructible} dense $T'$ de $T$, \struck{$\mathcal{U}'$ \ill{}
de $T \times_S T$ telle que $\varphi | \mathcal{U}'$ soit} et une
\emph{donnée de descente}
\[
\varphi' \colon X \times_S T' \xrightarrow{\ \sim\ } T' \times_S X
\]
telle que $\varphi = \varphi'$ \struck{\ill{}} au-dessus de
$\mathcal{U} \cap (T' \times_S T') = \mathcal{U}'$.
\note{énoncé encadré d'un trait vertical à gauche}

On peut supposer $\mathcal{U}$, $V$ \uncertain{ouverts} \ill{}
\emph{\uncertain{ou} \ill{} \uncertain{génériques}}.

Il existe une partie \uncertain{ouverte} dense $\mathcal{U}_1$ de
$\mathcal{U}$ telle que
\[
\tau \in \mathcal{U}_1 \Rightarrow V \cap p^{(3)-1}_{ij}(\tau)
\ \text{dense dans}\ p^{(3)-1}_{ij}(\tau)
\]
pour tout couple $(i,j)$ tel que $i \neq j$ [cf.\ \ill{} \ill{} \ill{}
\ill{} \ill{} \ill{}]. Remplaçant $\mathcal{U}$ par $\mathcal{U}_1$, on
\ill{} \ill{} \ill{} \ill{} que
$\tau \in \mathcal{U} \Rightarrow V \cap p^{(3)-1}_{ij}(\tau)$
\ill{} \ill{} \ill{} \uncertain{pour} \uncertain{tout} $(i,j)$ tel que
$i \neq j$. De même, \ill{} \ill{} \uncertain{partie}
\uncertain{constructible} $T'$ de $T$ telle que
\[
t \in T' \Rightarrow p_i^{-1}(t) \cap \mathcal{U}
\ \text{est dense dans}\ p_i^{-1}(t) \quad \text{pour } i = 1, 2 .
\]
Remplaçant $T$ par $T'$ (ce qui est \uncertain{loisible}) on peut
\uncertain{supposer} que
\[
t \in T \Rightarrow \mathcal{U} \cap p_i^{-1}(t) \ \text{dense dans}\ \ill{}
\quad \text{pour } i = 1, 2 .
\]
\struck{Enfin, on peut supposer \ill{} $\mathcal{U}$
\uncertain{symétrique}}

On désigne par $\xi$ un schéma de type fini variable sur $S$ [\ill{}
\ill{} \ill{} \ill{} : \ill{} \ill{} $\xi$]. On pose
\[
W = p^{(3)-1}_{12}(\mathcal{U}) \cap p^{(3)-1}_{23}(\mathcal{U})
\]
et \uncertain{on} \uncertain{considère}
\[
\pi \colon W \longrightarrow T \times T
\]
induit par $p_{13}$.

\page{27}

On considère $X_i = p_i^{*}(X)$ ($i = 1, 2$), et va définir un
isomorphisme
\[
\psi \colon \pi^{*}(X_1) \xrightarrow{\ \sim\ } \pi^{*}(X_2)
\]
compatible avec la \uncertain{projection} $\pi \colon W \to T \times_S T$,
i.e.\ \ill{} \ill{} \ill{} \ill{} $W \times_{\pi} W$ \ill{} \ill{}
\uncertain{isos}.

\struck{Pour cela, on \ill{} que}
\[
\pi^{*}(X_1) \simeq p^{(3)*}_{12}(X \times_S T | \mathcal{U}), \qquad
\pi^{*}(X_2) \simeq p^{(3)*}_{23}(T \times_S X | \mathcal{U})
\]
\note{ces deux formules, encadrées, sont barrées de deux traits obliques}

Pour cela, on considère les $p^{(3)*}_i(X)$ ($i = 1, 2, 3$), \ill{} on
\uncertain{note} que $p^{(3)*}_1(X) = \pi^{*}(X_1)$,
$p^{(3)*}_3(X) \simeq \pi^{*}(X_2)$, et \uncertain{prenons} \ill{} \ill{}
le composé
\[
p^{(3)*}_1(X) \xrightarrow{\ \sim\ } p^{(3)*}_2(X) \xrightarrow{\ \sim\ }
p^{(3)*}_3(X) .
\]
Pour prouver la compatibilité de $\psi$ avec $\pi$, \struck{on \ill{}
$W$} on \uncertain{écrit}
\[
W \times_{T \times_S T} W \subset T \times_S T \times_S T \times_S T ,
\]
\note{sous le produit, il écrit $t_1, t', t'', t_2$ pour les quatre
facteurs}
et on \uncertain{considère} le \struck{\uncertain{produit}}
\uncertain{sous-schéma} $W'$ de
$T \times_S T \times_S T \times_S T \times_S T$
(points $t_1, t', t'', t_2, t'''$), \struck{\ill{}} formé des
$(t_1, t', t'', t_2, t''')$ tels que
\[
(t''', t', t_1),\ (t''', t', t_2),\ (t''', t'', t_1),\ (t''', t'', t_2),\
(t''', t_1, t_2) \in V .
\]
\note{au-dessous, une première liste biffée : « $(t''', t_1), (t''', t_2),
(t''', t')$ \ill{} \ill{} $\mathcal{U}(\xi)$ » ; dans la liste retenue un
couple $(t''', t'')$ est aussi biffé, et $(t''', t_1, t_2)$ est ajouté en
indice à la fin}

On aura donc, si $t_1, t', t'', t_2, t'''$ sont comme il est spécifié
\uncertain{ci-dessus} : \struck{\ill{}}
\note{suit une formule biffée, commençant par
$\varphi_{t_2 t'} \varphi_{t' t_2} = (\varphi_{t_2 \ldots}$, que la page 28
reprend}

\begin{tikzcd}[column sep=small, row sep=small]
 & t' \arrow[dr, "\mathcal{U}"] & \\
t_1 \arrow[ur, "\mathcal{U}"] \arrow[r, "\mathcal{U}"] \arrow[dr, "\mathcal{U}"'] & t''' \arrow[u, "\mathcal{U}"'] \arrow[d, "\mathcal{U}"] \arrow[r, "\mathcal{U}"] & t_2 \\
 & t'' \arrow[ur, "\mathcal{U}"'] &
\end{tikzcd}
\note{schéma dans la marge gauche, en bas de la page ; la flèche de $t'''$
vers $t'$ est très repassée. Dans la marge gauche, en diagonale et
souligné : « \uncertain{Question} \ill{} \uncertain{rédaction} » ; plus
haut, des mots illisibles et « $p_i^{(3)}$ »}

\page{28}

\[
\varphi_{t_2 t'} \varphi_{t' t_1}
= \varphi_{t_2 t'} (\varphi_{t' t'''} \varphi_{t''' t_1})
= (\varphi_{t_2 t'} \varphi_{t' t'''}) \varphi_{t''' t_1}
= \varphi_{t_2 t'''} \varphi_{t''' t_1}
\]
On voit de même que
\[
\varphi_{t_2 t''} \varphi_{t'' t_1} = \varphi_{t_2 t'''} \varphi_{t''' t_1}
\]
d'où la compatibilité
\[
\varphi_{t_2 t'} \varphi_{t' t_1} = \varphi_{t_2 t''} \varphi_{t'' t_1} .
\]

On \uncertain{voit}, \uncertain{sauf} \uncertain{erreur}
(\ill{} \ill{}) que
$W' \to$ \struck{\ill{}} $W \times_{T \times_S T} W \simeq T^4$ \ill{}
\emph{\uncertain{surjectif}}.
\note{« $W \times_{T \times_S T} W \simeq T^4$ » est ajouté au-dessus de la
ligne, à la place d'un mot biffé}

Il s'ensuit, \uncertain{grâce} \ill{} \ill{} \ill{} \ill{} un morphisme de
$F$-descente, \ill{} que $\psi$ est bien compatible avec la
\uncertain{projection} $\pi \colon W \to T \times_S T$. On \uncertain{sait}
\uncertain{d'autre} \uncertain{part} que
\emph{$W \to T \times_S T$ est \uncertain{surjectif}}, \ill{}
\uncertain{plat}, il s'ensuit que c'est un morphisme de $F$-descente, donc
$\psi$ provient d'un \uncertain{isomorphisme}
\[
\varphi' \colon X_1 \xrightarrow{\ \sim\ } X_2 .
\]

Montrons que c'est une \struck{\ill{}} \uncertain{donnée} de descente,
\uncertain{donc} que
\[
\varphi'_{t_3 t_2} \varphi'_{t_2 t_1} = \varphi'_{t_3 t_1}
\quad (\text{sur } T \times_S T \times_S T) .
\]
\uncertain{Pour} \uncertain{cela}, on regarde la \struck{\uncertain{partie}}
\ill{} \ill{} \ill{} $W'''$ \struck{$W''$} de $T^6$ formée des
$t_1, t_2, t_3, t', t'', t'''$ tels que
\[
(t_1, t'),\ (t', t_2),\ (t_2, t'''),\ (t''', t_3) \in \mathcal{U},
\quad \text{et} \quad
(t_1, t', t'''),\ (t'', t''', t_3) \in V .
\]
On a $W''' \to T \times_S T \times_S T$, \ill{} \ill{} \ill{} \ill{} la
\uncertain{relation} \ill{} \ill{} : \ill{} \ill{} \ill{} \ill{} \ill{}
\ill{}.

\begin{tikzcd}[column sep=small, row sep=small]
 & & t''' \arrow[ddrr, bend left=40, "\mathcal{U}"] & & \\
 & t' \arrow[ur] \arrow[dr] & & t'' \arrow[ul, no head] \arrow[dr] & \\
t_1 \arrow[uurr, bend left=40, "\mathcal{U}"] \arrow[ur] & & t_2 \arrow[ur] & & t_3
\end{tikzcd}
\note{schéma dans la marge gauche ; le trait de $t''$ vers $t'''$ est
sans pointe}

\page{29}

\uncertain{Par} \uncertain{définition}, le \uncertain{premier}
\uncertain{membre} \uncertain{est} \ill{}
\[
(\varphi_{t_3 t''} \varphi_{t'' t_2})(\varphi_{t_2 t'} \varphi_{t' t_1})
\]
On \uncertain{a} :
\[
\varphi_{t'' t_2} \varphi_{t_2 t'} = \varphi_{t'' t'''} \varphi_{t''' t'}
\qquad [= \varphi'_{t'' t'}]
\]
\uncertain{donc} le \struck{\ill{}} \uncertain{premier}
\uncertain{membre} \ill{} \ill{} :
\[
(\varphi_{t_3 t''} \varphi_{t'' t'''})(\varphi_{t''' t'} \varphi_{t' t_1})
\]
La première parenthèse, \uncertain{puisque} $(t''', t'', t_3) \in V$,
\ill{} $\varphi_{t_3 t'''}$, et la \uncertain{deuxième} \uncertain{est}
\uncertain{égale} à $\varphi_{t''' t_1}$ ; d'où le \ill{} \uncertain{est}
$\varphi'_{t_3 t_1}$ \uncertain{par} \uncertain{définition}.
\note{petit schéma dans la marge gauche : $t'$ au-dessus de $t_1$ et $t_2$, relié à eux}

Enfin, \uncertain{remarquons} que \ill{} \ill{} \struck{\ill{}} \ill{}
\ill{} au-dessus de $\mathcal{U}$, $\varphi$ et $\varphi'$
\uncertain{coïncident} : \struck{\ill{}} \ill{} \ill{} \ill{} \ill{}
\ill{} \uncertain{que}
\[
V \xrightarrow{\ p_{13}\ } \mathcal{U}
\]
est surjectif, \struck{\uncertain{donc}} \ill{} \ill{} \ill{} \ill{}
\uncertain{ramené} à prouver l'\uncertain{assertion} \ill{} \ill{}
\uncertain{que} \ill{} \ill{} \ill{} $V$, \uncertain{laquelle} est
\uncertain{triviale}.

\uncertain{La} \uncertain{démonstration} \ill{} \ill{} \ill{}, \ill{}
\uncertain{facile}.

\emph{Condition \ill{} \uncertain{représentation}}. Supposons que pour
\uncertain{toute} partie \uncertain{constructible} \ill{} $T'$ de $T$,
\ill{} \ill{} $S'$ (\ill{}) de $S$, $T' \to S'$ \uncertain{soit}
\uncertain{un} morphisme de $F$-descente \emph{\uncertain{stricte}}.
Alors \uncertain{moyennant} \uncertain{remplacement} \ill{} $T'$ de $T$,
\ill{} \ill{}, on \uncertain{peut} \uncertain{trouver}
$Y \in F_{S'}$, \struck{\ill{}} \ill{} \uncertain{isom.}
\[
X \simeq Y \times_{S'} T' ,
\]
\uncertain{compatible} avec les \uncertain{données} de
\uncertain{descente}.
\note{« $\in F_{S'}$ » : le $S$ indicé est surchargé ; un ajout
interlinéaire à la fin de la page, sous la ligne, se lit
« \uncertain{descente} »}

\section*{Descente pour un morphisme fidèlement plat non fini}

\page{31}

\uncertain{Question} de la \uncertain{possibilité} de \uncertain{descente}
\uncertain{pour} \uncertain{un} \uncertain{morphisme} fid.\ plat
\emph{non fini}. \uncertain{Un} cas typique est le suivant :
\struck{\ill{}} $T$ \struck{\uncertain{plat}} \uncertain{étale}
\uncertain{gal.}\ fini sur $S$, de \uncertain{groupe} $G$, et
\struck{\ill{}} \uncertain{soit} $T'$ l'\uncertain{ouvert}
\struck{\ill{} $T$, \ill{}} $T - \Phi$, \uncertain{où}
\[
\Phi = f^{-1}(s) - (t) \qquad (f(t) = s)
\]
\ill{} $s$ \ill{} \uncertain{un} pt fermé de $S$.
\struck{\uncertain{On} \uncertain{a} \ill{} \ill{} $f^{-1}(s)$ \ill{}
$t^{c}$)} \uncertain{Alors} $T' \times_S T' \to$ \ill{}

\begin{tikzcd}[column sep=small, row sep=small]
 & X \arrow[d] \\
s \in S & T \arrow[l]
\end{tikzcd}
\note{schéma dans la marge gauche ; à gauche de $X$, un symbole
surchargé et une double flèche verticale, biffés}

On a
\[
T \times_S T = \Delta_T \times (G \times G)
\]
\note{entre $\Delta_T$ et $(G \times G)$, un signe de produit indicé
$G$, barré d'une croix ; le $T$ en indice de $\Delta$ est repassé}
\uncertain{formé} \ill{} \ill{} \ill{} \ill{} \ill{} \ill{}
\[
\mathcal{U}_g = (g \times e)\, \Delta_T ,
\]
\struck{\ill{}} \ill{} \ill{} \ill{} \uncertain{isomorphe}
\uncertain{à} $T$ \uncertain{par} \ill{}
\struck{\ill{} \ill{} $T \to \Delta_T \xrightarrow{g \times e}$}
\note{la fin de la phrase est remplacée par un ajout au-dessus de la
ligne, en partie illisible}
\uncertain{montrant}

\begin{tikzcd}
T \arrow[r, "\varphi_g"] \arrow[d, "\Delta"'] & \mathcal{U}_g \arrow[d, "\mathrm{ind.}"] \\
T \times_S T \arrow[r, "g \times e"'] & T \times_S T
\end{tikzcd}

et \uncertain{pour} \ill{} les \ill{} \uncertain{correspondants}
\uncertain{aux} \uncertain{projections} de $\varphi_g$ \ill{}
\ill{} \uncertain{respectivement} $g$ et l'identité, \ill{} \ill{}
\[
\varphi_g^{-1}\bigl(\mathcal{U}_g \cap (T' \times_S T')\bigr)
= g^{-1}(T') \cap T' .
\]
\uncertain{Il} \uncertain{est} \ill{} \ill{} $T'$ si $g = e$,
\uncertain{et} $T - f^{-1}(s) = T''$ si $g \neq e$.
\emph{Ainsi}, $T' \times_S T'$ \emph{est}
\struck{\uncertain{réunion}} \uncertain{somme} \ill{}
\emph{\ill{} \uncertain{sa} \uncertain{diagonale}} $\simeq T$ ; et
\emph{d'\uncertain{ouverts}}
$\mathcal{U}'_g$ ($g \neq e$) \emph{\ill{} \uncertain{isomorphes}} :
$T'' = T - f^{-1}(s)$.
\uncertain{Une} \uncertain{donnée} de descente \uncertain{sur}
\uncertain{un} $X / T'$ \uncertain{relativement} \ill{} : \ill{}
\ill{} \ill{}
\note{à la fin de la page le tracé se lit « $X / T'$ » ; « $T''$ » est
souligné}

\page{32}

\struck{un \ill{} \uncertain{isomorph.}\ d'\uncertain{isomorphismes}}
\ill{} \uncertain{isom.}
\[
X \times_S T' \simeq T' \times_S X
\]
i.e.\ un \uncertain{syst.}\ d'isomorphismes des deux
\uncertain{images} \uncertain{inverses} de $X$ sur les
$\mathcal{U}'_g$ ; sur la \uncertain{diagonale} $\mathcal{U}'_e$, ce
\uncertain{doit} être l'identité, et pour $g \neq e$, \uncertain{ce}
doit être un isomorphisme
\[
X | T'' \simeq g^{*}(X | T'') .
\]
Il \uncertain{faut} écrire la condition \ill{} \ill{} \ill{}
\uncertain{transitivité}. \ill{} \ill{}
\[
T \times_S T \times_S T \simeq \Delta_3 T \times (G \times G \times G) .
\]
\note{ici encore le signe de produit indicé $G$ devant
$(G \times G \times G)$ est barré d'une croix}
\uncertain{Soit}
\[
V_{g, g'} = (g \times g' \times e) \Delta_3 T ,
\]
il s'identifie par $(g \times g' \times e) \circ \Delta$ avec $T$,
\uncertain{les} \ill{} \ill{} \ill{} \ill{} \ill{} \ill{} $g, g', e$ ;
\ill{}
\[
V'_{g, g'} = V_{g, g'} \cap (T' \times_S T' \times_S T')
\]
s'identifie \uncertain{à}
\[
g^{-1}(T') \cap g'^{-1}(T') \cap T' ,
\]
\ill{} \ill{} $T''$ \uncertain{sauf} \uncertain{pour} $g = g' = e$,
\ill{} \ill{} \ill{} la diagonale $\simeq T'$.
$T' \times_S T' \times_S T'$ \uncertain{est} \ill{} \ill{} \ill{}
$V'_{g,g'}$, \ill{} \ill{} \ill{} \uncertain{isomorphes} \uncertain{à}
$T''$. On \uncertain{constate} \ill{} que la
\emph{transitivité},

\page{33}

\uncertain{vérification} \ill{} \uncertain{pour} \uncertain{deux}
\uncertain{morphismes} \ill{} \uncertain{précédemment} \uncertain{sur}
$T' \times_S T' \times_S T'$ \uncertain{soient} \uncertain{identiques},
\uncertain{étant} \ill{} \uncertain{acquis} \uncertain{sur} la
\uncertain{diagonale} $V'_{e,e}$, \ill{} il \uncertain{suffit} à
l'exprimer \uncertain{sur} les \uncertain{autres}
\uncertain{composantes}. Mais \uncertain{alors} il ne \uncertain{porte}
que sur les $T''$, \uncertain{et} \struck{la \uncertain{superposition}
\ill{}} \ill{} \ill{} \ill{} \struck{\ill{}} une \uncertain{donnée} de
\uncertain{descente} \struck{\ill{} \ill{}} \ill{} sur $X | T''$,
\ill{} \ill{} \uncertain{une} \uncertain{donnée} de descente, i.e.\
\uncertain{les} \uncertain{isom.}
\[
g^{-1}(X) \simeq X
\]
\uncertain{proviennent} d'\uncertain{opérations} de $G$ sur $X | T''$,
compatibles avec \uncertain{ses} \uncertain{opérations}
\struck{\uncertain{sur}} \struck{$X$} $T''$. \emph{Ainsi, une
donnée de descente sur $X$ pour $f \colon T' \to S$ \uncertain{est} la
même chose qu'une donnée de descente sur $X | T''$, i.e.\ que \struck{\ill{}}
la donnée d'\uncertain{opérations} de $G$ sur $X | T''$, compatibles avec
ses \uncertain{opérations} sur $T''$}. Si p.ex.\ $X$ est
\uncertain{proj.}\ sur $T'$, \uncertain{cela} \ill{} \ill{} \ill{}
\uncertain{que} la \uncertain{donnée} d'un \uncertain{isom.}
\[
X | T'' \simeq X_0 \times_{S''} T''
\]
où $S'' = S - (s)$, $X_0$ \uncertain{étant} un schéma
\uncertain{proj.}\ sur $S''$.

\page{34}

La \uncertain{question} d'effectivité de la \uncertain{descente}
\uncertain{est} \uncertain{donc} \uncertain{équivalente} à la
\uncertain{suivante} :

\uncertain{On} \uncertain{a} $X'_0$ sur $S - s$, on \uncertain{étend}
\[
X'_0 \times_{S - s} \bigl(T - f^{-1}(s)\bigr)
\]
en un point $t$ de $f^{-1}(s)$ de façon à \uncertain{obtenir} $X$ sur
\[
T' = \bigl(T - f^{-1}(s)\bigr) \cup \{t\} ,
\]
\uncertain{par} \uncertain{exemple} \uncertain{sur} $T$ !
\uncertain{Montrer} que $X$ \uncertain{est} \uncertain{isomorphe} à
l'image \uncertain{réciproque} d'un $X_0$ sur $S$ \uncertain{prolongeant}
$X_0$.
\note{le dernier $X_0$ est à lire, semble-t-il, comme le $X'_0$ du début
du paragraphe}

\end{document}
